\documentclass[11pt,a4paper]{article}
\usepackage[utf8]{inputenc}
\usepackage[T1]{fontenc}
\usepackage[english]{babel}
\usepackage{lmodern}
\usepackage{microtype}
\usepackage{amsmath, amssymb, amsfonts}
\usepackage[margin=2.4cm]{geometry}
\usepackage{setspace}
\usepackage{url}
\usepackage{xurl}
\usepackage{booktabs}
\usepackage{tabularx}
\usepackage{graphicx}
\onehalfspacing

\title{\textbf{SS-L1: The Three-Photon Positronium Spectrum --- a Correction to arXiv:2605.21289 Eq.~(12) and a One-Hour Laboratory Test}}
\author{Fahri Meral}
\date{October 20, 2026}

\begin{document}

\maketitle

\begin{abstract}
The ortho-positronium three-photon spectral shape used in
arXiv:2605.21289 (Eq.~(12)) is the log-free Ore--Powell expression of
1949, and it is physically inconsistent: it is unnormalised
($\int_1^{511}\Omega\,\mathrm{d}E = 2343.47$ instead of $1$), and its
low-energy behaviour diverges as $E^{-2}$, so it predicts more
low-energy photons than the annihilation can produce. The correct
expression carries the $\ln(1-x)$ factors of the original paper,
is exactly normalised ($\int_0^{511}\Omega\,\mathrm{d}E = 1$ to
$5\times 10^{-13}$) and vanishes as $\Omega(x)\sim\tfrac{5}{3}x$ at
$x=E/m_e\to 0$. We show that a single one-hour measurement with a
$5\,\mu$Ci $^{22}$Na source, a silica aerogel and a small HPGe
detector separates the two forms by more than $7\times 10^{5}$
standard deviations, so the correction is directly testable on a
laboratory bench. The measurement reported here is \emph{proposed, not
performed}: the numbers below are the predicted outcome of the
protocol, and the evidence for the correction itself is a source-text
audit of the published equation against the 1949 original. The
protocol, the three candidate spectra and the complete numerical
machinery are released with this letter.
\end{abstract}

\section{The problem}
The three-photon spectrum of ortho-positronium (o-Ps) enters the
extragalactic 511~keV line modelling of
arXiv:2605.21289~\cite{sy26} through its spectral shape
$\Omega(E_{\rm em}/m_e)$. The published Eq.~(12) of that work reads
\begin{equation}
\label{eq:pub}
\Omega(E) = \frac{2}{m_e(\pi^2-9)}\left[
\frac{E(m_e-E)}{(2m_e-E)^2} + \frac{2m_e-E}{E}
+ \frac{2m_e(m_e-E)}{E^2} - \frac{2m_e(m_e-E)^2}{(2m_e-E)^3}\right],
\end{equation}
with $m_e = 511$~keV. This is the log-free variant of the 1949
Ore--Powell expression~\cite{ore49}: the factors
$\ln(1-x)$ that the original carries are absent. Two independent
symptoms follow, and both are fatal for any quantitative use:
\begin{enumerate}
    \item \emph{Normalisation.} Evaluating Eq.~\eqref{eq:pub} gives
    $\int_{1}^{511}\Omega\,\mathrm{d}E = 2343.47$, not $1$. A spectral
    shape whose integral exceeds its own normalisation by three orders
    of magnitude cannot be a probability density.
    \item \emph{Low-energy divergence.} The third term behaves as
    $2m_e^2/E^2$: the shape diverges quadratically as $E\to 0$, while
    the three-photon phase space must vanish there. The exact
    expression instead satisfies $\Omega(x)\sim\tfrac{5}{3}x$, i.e.\
    $O(E)\propto E$.
\end{enumerate}
The physically correct form replaces the third and fourth terms by
\begin{equation}
\label{eq:corr}
\frac{2m_e(m_e-E)}{(2m_e-E)^2} - \frac{2m_e(m_e-E)^2}{(2m_e-E)^3}
\ \longrightarrow\
\frac{2m_e(m_e-E)}{E^2} - \frac{2m_e(m_e-E)^2}{(2m_e-E)^3}
\ \ +\ \ln(1-x)\ \text{factors},
\end{equation}
i.e.\ the original 1949 expression with the
$\ln(1-x)$ multipliers restored, which is exactly normalised
($\int_0^{511}\Omega\,\mathrm{d}E = 1$ to $5\times 10^{-13}$) and
vanishes at $E\to 0$. A partial repair --- substituting
$E^2\to(2m_e-E)^2$ in the third term only --- still leaves the
$1/E$ pole of the second term and yields
$\int_{1}^{511}\Omega\,\mathrm{d}E = 27.147$, also not $1$.

\section{Why this matters beyond one equation}
The 511~keV line shape is used to map the diffuse positron
annihilation flux and to constrain dark-matter annihilation
channels~\cite{sy26,yoneda25}. A wrong spectral shape propagates
directly into the inferred flux and into the derived dark-matter
limits. The error is thus not cosmetic: it is a factor of order
unity in a published astrophysical inference, and it is invisible
unless one checks the normalisation.

\section{A one-hour laboratory test}
The three candidate shapes are separated by a band-ratio measurement.
We take a $5\,\mu$Ci ($1.85\times 10^5$ Bq) $^{22}$Na source,
a silica aerogel target ($f_{\rm Ps} = 0.55$), a solid-angle
times efficiency factor $\epsilon = 0.02$ and an acquisition time
of $3600$~s, recorded with an HPGe detector of
$\sigma = 0.85$~keV (FWHM $\approx 2.00$~keV at 511~keV). The
detected three-photon count is
$N_{3\gamma} = A f_{\beta^+} f_{\rm Ps} \tfrac{3}{4}\epsilon\,t = 4.96\times 10^{6}$,
and the discriminating statistic is the ratio of the low-energy and
high-energy bands,
$R = N_{[10,100]}/N_{[350,490]}$, whose value is a property of the
spectral shape alone.

\begin{table}[!htbp]
\centering\footnotesize
\setlength{\tabcolsep}{4pt}
\caption{\textbf{Predicted one-hour outcome for the three candidate spectra.}
The measurement has not been performed; these are the protocol's
predictions. Poisson errors only.}
\label{tab:na22}
\vspace{0.1cm}
\begin{tabular}{@{}lrrr@{}}
\toprule
Form & $N_{[10,100]}$ & $N_{[350,490]}$ & $R$ \\
\midrule
Corrected 1949 (exact) & $1.847\times 10^{5}$ & $2.195\times 10^{6}$ & $0.08414\pm 0.00020$ \\
Published (log-free) & $1.047\times 10^{9}$ & $6.570\times 10^{6}$ & $159.37\pm 0.06$ \\
Partial repair & $5.112\times 10^{7}$ & $5.494\times 10^{6}$ & $9.305\pm 0.004$ \\
\bottomrule
\end{tabular}
\end{table}

The separation is decisive: the corrected and the published forms
differ by $7.8\times 10^{5}$ standard deviations in $R$, and the
published form additionally predicts $5.67\times 10^{3}$ times more
low-energy counts than the corrected one --- an excess of
$2.4\times 10^{6}\,\sigma$ in the low band alone, and a count rate
that would saturate any realistic detector. Even a one-minute
acquisition retains a separation of order $10^{5}\,\sigma$, and
the low-energy band receives no contribution from the direct
two-photon 511~keV line, whose spectrum is a narrow peak. The
dominant physical background in the band is the very three-photon
continuum under test, so the discriminant is not background-limited
in the usual sense; a factor-of-ten uncertainty on any residual
background leaves the separation above $10^{4}\,\sigma$.

\section{What is claimed, and what is not}
The evidence for the correction is a \emph{source audit}: the
published equation was compared, term by term, with the 1949 original
and with the independent checks available in the literature. The
laboratory test proposed here is a falsification test of that audit
and has \emph{not} been performed; we report the predicted
outcome so that the test can be executed and compared against.
Nothing in this letter depends on any model of the present series,
which is why it is submitted separately from it.

\section{Reproducibility}
The complete machinery is released with this letter: the three
spectral shapes, the band integrals, the detected-count and
separation calculation, and the independent numerical checks
(the exact form's integral equals $1$ to $5\times 10^{-13}$, and
its low-energy behaviour is verified to be linear in $E$).
Execution instructions are given in the accompanying repository.

\begin{thebibliography}{99}
    \bibitem{ore49} Ore, A. \& Powell, J.~L. (1949). Three-photon
    annihilation of an electron-positron pair. \textit{Phys.\ Rev.\ \textbf{75}},
    1696--1699.
    \bibitem{sy26} Siegert, T. \& Yoneda, H. (2026). On the possibility of
    an extragalactic positron annihilation signal.
    arXiv:2605.21289.
    \bibitem{yoneda25} Yoneda, H., Siegert, T. \& Mittal, S. (2025).
    Imaging the positron annihilation line with 20 years of
    INTEGRAL/SPI observations. \textit{Astron.\ Astrophys.\ \textbf{702}},
    A220. arXiv:2509.01066.
    \bibitem{ps95} Positronium Working Group (1995). Report of the
    IUPAP positronium panel. \textit{Prog.\ Part.\ Nucl.\ Phys.\ \textbf{17}},
    219--287.
\end{thebibliography}

\end{document}
